\chapter{Implementation} \label{app:implementation}

\section{Function Block Files}

\begin{lstlisting}[language=XML, caption={Example \textit{xml} definition of a function block}, label={lst:fb-xml}]
<FBType Name="NEW_FB_1">
  <InterfaceList>
    <EventInputs>
      <Event Name="INIT" Type="Event"/>
      <Event Name="READ" Type="Event">
        <With Var="VARIABLE_1"/>
        <With var="VARIABLE_2"/>
        <With var="VARIABLE_3"/>
      </Event>
    </EventInputs>
    <EventOutputs>
      <Event Name="INIT_O" Type="Event"/>
      <Event Name="READ_O" Type="Event">
        <With Var="DATA_1"/>
      </Event>
    </EventOutputs>
    <InputVars>
      <VarDeclaration Name="VARIABLE_1" Type="STRING"/>
      <VarDeclaration Name="VARIABLE_2" Type="REAL"/>
      <VarDeclaration Name="VARIABLE_3" Type="INT"/>
    </InputVars>
    <OutputVars>
      <VarDeclaration Name="DATA_1" Type="STRING"/>
    </OutputVars>
  </InterfaceList>
</FBType>
\end{lstlisting}


\begin{lstlisting}[language=Python, caption={Example \textit{Python} code of a function block}, label={lst:fb-py}]
class NEW_FB_1:
  def __init__(self):
    pass

  def schedule(self, event_name, event_value, variable_1, variable_2, variable_3):
    if event_name == "INIT":
      return [event_value, None, ""]

    elif event_name == "READ":
      return [None, event_value, ""]
\end{lstlisting}

\section{Synchronisation}

\begin{lstlisting}[language=JSON, caption={Sync configuration}, label={lst:sync-config}]
{
  master-fbs-path: 'sync/fbs',
  dinasores: [
    {
      address: 'x.x.x.x',
      port: 22,
      username: 'root',
      password: 'dinasore',
      dinasore-path: '/root/dinasore'
    },
    {
      address: 'y.y.y.y',
      port: 22,
      username: 'root',
      password: 'dinasore',
      dinasore-path: '/root/dinasore'
    }
  ]
}
\end{lstlisting}

\section{Communication} \label{app:implementation:communication}

\begin{lstlisting}[language=Python, caption={build\_message() in Python}, label={lst:build-message-py}]
def build_message(message_payload, configuration_name):
  # build the first part of the header
  hex_input = '{:04x}'.format(len(configuration_name))
  second_byte = int(hex_input[0:2], 16)
  third_byte = int(hex_input[2:4], 16)
  response_len = struct.pack('BB', second_byte, third_byte)
  response_header_1 = b''.join([b'\x50', response_len])

  # build the second part of the header
  hex_input = '{:04x}'.format(len(message_payload))
  second_byte = int(hex_input[0:2], 16)
  third_byte = int(hex_input[2:4], 16)
  response_len = struct.pack('BB', second_byte, third_byte)
  response_header_2 = b''.join([b'\x50', response_len])

  # join the 3 parts
  response = b''.join([response_header_1, configuration_name, response_header_2, message_payload])
  return response
\end{lstlisting}

\begin{lstlisting}[language=JavaScript, caption={build\_message() in JavaScript}, label={lst:build-message-js}]
function build_message(message_payload, configuration_name) {
  const configBuffer = Buffer.from(configuration_name, 'utf-8');
  const messageBuffer = Buffer.from(message_payload, 'utf-8');

  // Header 1: [0x50][len_hi][len_lo]
  const header1 = Buffer.alloc(3);
  header1[0] = 0x50;
  header1.writeUInt16BE(configBuffer.length, 1);

  // Header 2: [0x50][len_hi][len_lo]
  const header2 = Buffer.alloc(3);
  header2[0] = 0x50;
  header2.writeUInt16BE(messageBuffer.length, 1);

  return Buffer.concat([
    header1,
    configBuffer,
    header2,
    messageBuffer
  ]);
}
\end{lstlisting}

\section{Send Messages} \label{app:send-messages}
First we need to send the two commands in \Cref{lst:deploy-setup} to every device. A \texttt{QUERY} request and a \texttt{CREATE} request that will create a \textit{EMB\_RES} resource which will be automatically connected to the \textit{INIT} port of all other function block nodes to initialize them. Both these commands will have an empty configuration name, as \textit{EMB\_RES} has not been created yet and therefore cannot be bounded to.

\begin{lstlisting}[language=XML, caption={Setup deployment commands}, label={lst:deploy-setup}]
<Request Action="QUERY" ID="0">
  <FB Name="*" Type="*"/>
</Request>

<Request Action="CREATE" ID="1">
  <FB Name="EMB_RES" Type="EMB_RES"/>
</Request>
\end{lstlisting}

Next, depending on which commands are mapped to each machine, we \texttt{CREATE} all the mapped nodes and \texttt{WRITE} the initial input values with the commands in \Cref{lst:deploy-create-nodes}. The configuration name in these and all following commands will be "\textit{EMB\_RES}" as they need to be bound to it in order to be initialised and ran.

\begin{lstlisting}[language=XML, caption={\texttt{CREATE} nodes deployment commands}, label={lst:deploy-create-nodes}]
<Request Action="CREATE" ID="2">
  <FB Name="SENSOR_SIMULATOR" Type="SENSOR_SIMULATOR"/>
</Request>

<Request Action="WRITE" ID="3">
  <Connection Destination="SENSOR_SIMULATOR.OFFSET" Source="12"/>
</Request>

<Request Action="CREATE" ID="4">
  <FB Name="MOVING_AVERAGE" Type="MOVING_AVERAGE"/>
</Request>

<Request Action="WRITE" ID="5">
  <Connection Destination="MOVING_AVERAGE.WINDOW" Source="5"/>
</Request>
\end{lstlisting}

Then the connections are created between the mapped nodes. A \texttt{Source} a \texttt{Destination} must be provided, with the format \texttt{"NODE\_NAME.SLOT\_NAME"}, using the commands in \Cref{lst:deploy-create-links}.

\begin{lstlisting}[language=XML, caption={\texttt{CREATE} links deployment commands}, label={lst:deploy-create-links}]
<Request Action="CREATE" ID="6">
  <Connection Destination="SENSOR_SIMULATOR.READ_0" Source="MOVING_AVERAGE.RUN"/>
</Request>

<Request Action="CREATE" ID="7">
  <Connection Destination="SENSOR_SIMULATOR.VALUE" Source="MOVING_AVERAGE.VALUE"/>
</Request>
\end{lstlisting}

Finally the \texttt{START} command in \Cref{lst:deploy-start} can be issued, which will request for the start of the program execution. For this command the configuration name used will also be "\textit{EMB\_RES}" as it is this resource that will begin the execution of the program.

\begin{lstlisting}[language=XML, caption={\texttt{START} deployment command}, label={lst:deploy-start}]
<Request Action="START" ID="8"/>
\end{lstlisting}

After each request is sent, a response will be returned by the \textit{DINASORE}. It is important for the \ac{IDE} to listen for these responses and only send the next request once the response corresponding to the previous request is received. The responses have the format shown in \Cref{lst:deploy-response}. The \texttt{ID} of the response will be the same of the associated request. The format of the responses is the same as the requests and the configuration name that is used is empty.

\begin{lstlisting}[language=XML, caption={\texttt{Response} after a deployment command is sent}, label={lst:deploy-response}]
<Response ID="8"/>
\end{lstlisting}

\section{Watch} \label{app:watch}
Note in \Cref{lst:deploy-create-watch} that the \texttt{ID} is reset after deploying.

\begin{lstlisting}[language=XML, caption={\texttt{CREATE} watches after deploying}, label={lst:deploy-create-watch}]
<Request ID="0" Action="CREATE">
  <Watch Source="SENSOR_SIMULATOR.VALUE" Destination="*" />
</Request>

<Request ID="1" Action="CREATE">
  <Watch Source="MOVING_AVERAGE.VALUE_MA" Destination="*" />
</Request>
\end{lstlisting}

After creating the watches we can send a request to receive the values using the command in \Cref{lst:deploy-read-watch}. This will return the response in \Cref{lst:deploy-watch}, which by parsing the values, allows the prototype to display them in the \textit{Graph View}.

\begin{lstlisting}[language=XML, caption={Request to \texttt{READ} watches}, label={lst:deploy-read-watch}]
<Request ID="2" Action="READ">
  <Watches/>
</Request>
\end{lstlisting}

\begin{lstlisting}[language=XML, caption={Response to \texttt{READ} watches}, label={lst:deploy-watch}]
<Response ID="2">
	<Watches>
		<Resource name="EMB_RES">
			<FB name="SENSOR_SIMULATOR">
				<Port name="VALUE">
					<Data value="13.75" />
				</Port>
			</FB>
			<FB name="MOVING_AVERAGE">
				<Port name="VALUE">
					<Data value="13.3" />
				</Port>
			</FB>
		</Resource>
	</Watches>
</Response>
\end{lstlisting}
